\section[\texorpdfstring{Approximation by the wavelet spaces $\WJ$}{Approximation by the wavelet spaces W\_J}]{Approximation by the wavelet \texorpdfstring{spaces $\WJ$}{spaces W\_J}}
\setcounter{equation}{0}

\subsection*{Proof of Lemma 2}

Fix a pair $(\cv,\md)$ and write $g = \gcomp{\cv}{\md}$ and
$\theta_{\lev\tsl} = \coefs{\cv}{\md}{\lev}{\tsl}$.

\emph{Tail sum.} By Step 2 below the coefficients are square summable, so
$g \in L_2[0,1]$ and, the system being complete and $g$ orthogonal to the
constant, $g = \sum_{\lev\ge0}\sum_{\tsl}\theta_{\lev\tsl}\wav{\lev}{\tsl}$ in
$L_2$. The orthogonal projection onto $\WJ$ retains the levels below $J$, and
Parseval's identity gives
$\norm{g-\PiJ g}^2_{L_2[0,1]} = \sum_{\lev\ge J}
\norm{\theta_{\lev\cdot}}_2^2$.

\emph{From $\ell_\pi$ to $\ell_2$ inside a level.} In $\R^{2^{\lev}}$ one has
$\norm{a}_2 \le \norm{a}_{\pi}$ when $\pi \le 2$, and
$\norm{a}_2 \le (2^{\lev})^{1/2-1/\pi}\norm{a}_{\pi}$ by Hölder's inequality
when $\pi > 2$. In both cases $\norm{\theta_{\lev\cdot}}_2 \le
2^{\lev(1/2-1/\pi)_{+}}\norm{\theta_{\lev\cdot}}_{\pi}$, and
Assumption 5 gives
\[
  \norm{\theta_{\lev\cdot}}_2 \;\le\;
  C_g\,2^{-\lev\{s+1/2-1/\pi-(1/2-1/\pi)_{+}\}}
  \;=\; C_g\,2^{-\lev\seff},
\]
since the exponent equals $s$ when $\pi \ge 2$ and $s-(1/\pi-1/2)$ when
$\pi < 2$. This is where the loss $(1/\pi-1/2)_{+}$ appears.

\emph{Geometric sum and density.}
$\sum_{\lev\ge J}C_g^2 2^{-2\lev\seff} = C_g^2 2^{-2J\seff}/(1-2^{-2\seff})$,
finite because $\seff > 0$, which is the first display of Lemma 2. For the
second, write $f-f_J = \sum_{\cv,\md}X_{\cv}h_{\cv\md}(U_{\md})$ with
$h_{\cv\md} = \gcomp{\cv}{\md}-\PiJ\gcomp{\cv}{\md}$, the levels cancelling.
By the Cauchy-Schwarz inequality over the $pq$ terms and
$|X_{\cv}| \le B_X$,
\[
  \E\{(f-f_J)^2\} \;\le\; pq\sum_{\cv,\md}
  \E\{X_{\cv}^2h_{\cv\md}^2(U_{\md})\}
  \;\le\; pq\,B_X^2\sum_{\cv,\md}\E\{h_{\cv\md}^2(U_{\md})\},
\]
and $\E\{h^2(U_{\md})\} = \int_0^1 h^2 p_{U_{\md}} \le C_U\norm{h}^2_{L_2[0,1]}$,
the marginal density being bounded by $C_U$ because integrating the joint
density over the other $q-1$ coordinates of $[0,1]^{q-1}$ does not increase
the bound. One factor $pq$ comes from the Cauchy-Schwarz inequality and the
other from the sum. \qed

The uniform version is used in Section S5 to bound the empirical norm of the
bias without an appeal to Markov's inequality.

\begin{lemma}[Uniform approximation error]
\label{lem:sup}
Under Assumptions 4 and 5 with $s > 1/\pi$, each component has a continuous
version and, for every $J \ge 0$,
\[
  \norm{\gcomp{\cv}{\md}-\PiJ\gcomp{\cv}{\md}}_{\infty}
  \;\le\; \frac{(L-1)\norm{\psi}_{\infty}C_g}{1-2^{-(s-1/\pi)}}\,
  2^{-J(s-1/\pi)} ,
\]
and therefore
$\normn{\bb} \le \norm{f-f_J}_{\infty} \le pq B_X(L-1)\norm{\psi}_{\infty}
C_g\,2^{-J(s-1/\pi)}/\{1-2^{-(s-1/\pi)}\}$ for every sample.
\end{lemma}

\begin{proof}
At a point $u$ and a level $\lev$ at most $L-1$ periodized translates are
nonzero, because $\psi(2^{\lev}u-\tsl)\ne0$ requires
$2^{\lev}u-\tsl \in{}$ $(1-L/2,\,L/2)$, so
$\sum_{\tsl}|\wav{\lev}{\tsl}(u)| \le (L-1)2^{\lev/2}\norm{\psi}_{\infty}$.
With $\max_{\tsl}|\theta_{\lev\tsl}| \le \norm{\theta_{\lev\cdot}}_{\pi} \le
C_g2^{-\lev(s+1/2-1/\pi)}$ the level contributes at most
$(L-1)\norm{\psi}_{\infty}C_g2^{-\lev(s-1/\pi)}$ in the supremum norm. Since
$s > 1/\pi$ the series converges absolutely and uniformly, so its sum is
continuous and agrees almost everywhere with $g$; the tail from level $J$ on
is the geometric sum of the displayed bounds, and the statement for $f-f_J$ is
the triangle inequality over the $pq$ terms.
\end{proof}

The next proposition is the one referred to in Section 2.2 of the paper: it
quantifies what the periodized basis charges for a component that does not
match at the two ends of the interval.

\begin{proposition}[Cost of periodization]
\label{prop:periodization}
Let the basis be as in Assumption 4 and let
$g \in \Besov{s}{\pi}{r}[0,1]$, the Besov space of the interval and not of the
circle, with $s > 1/\pi$, $N > s$ and $\int_0^1 g = 0$. Then there is a
constant $C$ depending on $g$, on $\psi$ and on $(s,\pi)$ such that, for every
$J$ with $2^J \ge L-1$,
\[
  \norm{g-\PiJ g}^2_{L_2[0,1]} \;\le\; C\,2^{-2J\seff}
  \;+\; 2(L-1)\norm{\psi}_1^2\norm{g}^2_{\infty}\,2^{-J},
\]
so that the $L_2$ rate is $2^{-J\min(\seff,1/2)}$. If $g$ is continuous with
$g(0) \ne g(1)$, the second term is of the exact order $2^{-J}$ and
$\norm{g-\PiJ g}_{\infty}$ does not converge to zero.
\end{proposition}

\begin{proof}
Let $\tilde g$ be the one periodic extension of $g$, so that
$\theta_{\lev\tsl} = \int_{\R}\tilde g(y)2^{\lev/2}\psi(2^{\lev}y-\tsl)\,dy$.
At a level with $2^{\lev} \ge L-1$, separate the interior translates, whose
support lies inside $[0,1]$, from the at most $L-1$ boundary translates that
cross the identification of $0$ with $1$. For an interior translate,
$\theta_{\lev\tsl} = \inner{Eg}{\psi_{\lev\tsl}}_{\R}$ for any extension $Eg$
of $g$ to the line; taking $E$ a bounded extension operator from
$\Besov{s}{\pi}{r}[0,1]$ to $\Besov{s}{\pi}{r}(\R)$
\citep[Chapter 3]{Triebel-1983} and using the wavelet characterization of the
latter, which for the implication from the function to the sequence requires
only $N > s$ vanishing moments, gives
$\norm{(\theta_{\lev\tsl})_{\rm int}}_{\pi} \le C_12^{-\lev(s+1/2-1/\pi)}$,
whose tail sum is bounded as in Lemma 2. For a boundary translate,
$s > 1/\pi$ makes $g$ bounded and
$|\theta_{\lev\tsl}| \le \norm{g}_{\infty}\norm{\psi_{\lev\tsl}}_1 =
\norm{g}_{\infty}\norm{\psi}_12^{-\lev/2}$, and there are at most $L-1$ of
them per level, so their tail sum is at most
$2(L-1)\norm{\psi}_1^2\norm{g}_{\infty}^22^{-J}$. For the exact order, if $g$
is continuous with $g(0) \ne g(1)$ then on the support of a boundary translate
$\tilde g$ equals $g(1)+o(1)$ to the left of the identification point and
$g(0)+o(1)$ to the right, so
$\theta_{\lev\tsl} = 2^{-\lev/2}\{g(0)-g(1)\}\int_{y>a_{\tsl}}\psi(y)\,dy +
o(2^{-\lev/2})$; since $\int\psi = 0$ and $\psi$ vanishes on no subinterval of
its support, that partial integral is nonzero for some translate, and the
boundary contribution of the level is of order $2^{-\lev}$. The failure of
uniform convergence is that of any wavelet series at a jump.
\end{proof}

Under Assumption 5 nothing changes, since that assumption is already a
periodic one; what Proposition S3.1 gives is the price of its failure. With
the basis of \citet*{cohen1993wavelets} the space $\VJ$ reproduces polynomials
of degree below $N$ and its basis characterizes $\Besov{s}{\pi}{r}[0,1]$
without periodicity \citep[Section 3.9]{Cohen-2003}, so Lemma 2 holds with
that space and the same proof.

%%%%%%%%%%%%%%%%%%%%%%%%%%%%%%%%%%%%%%%%%%%%%%%%%%%%%%%%%%%%%%%%%%%%%%%%%%%%%%
% derivations/03-desenho-produtos.tex
